% !PRINTED none
% part-title page (no running head, no printed page number): centred heading
% "Nachträge." with a horizontal rule beneath it, then subsection "I." and
% "Irreducibilität der Kreistheilungsgleichung.", followed by body text.
\begin{center}
{\Large \gesperrt{Supplements.}}

\rule{4em}{0.4pt}

I.

\textbf{\gesperrt{Irreducibility of the Cyclotomic Equation.}}
\end{center}

\S.\,134 of the first volume stands in need, in its last part, which treats of the irreducibility of the cyclotomic equation in the general case in which the order of the root of unity is a number composed of several distinct primes, of a correction, which shall be supplied here; and we ask the reader to complete the passage named in accordance with what is set out here.

If $n=ab$ is put, $a$ and $b$ being relatively prime, and if $\varrho, \alpha, \beta$ are primitive $n$th, $a$th, $b$th roots of unity, then $\varrho=\alpha\beta$, and it is proved in the passage cited that, if $\varrho$ is a root of a rational equation $\Phi(x)=0$, then among the roots of this equation $\varrho=\alpha\beta$ all primitive $a$th roots of unity $\alpha$ and all primitive $b$th roots of unity $\beta$ must occur. But it does not follow from this, as was there erroneously assumed, that among them \gesperrt{all} primitive $n$th roots of unity occur, because it has not been shown that every $\alpha$ occurs combined with every $\beta$.

To fill this gap, let an entirely elementary proof first find a place here.\footnote{\gesperrt{Arndt}, \gesperrt{Crelle}'s Journ., Bd. 56, S. 178 (1859). \gesperrt{Lebesgue}, \gesperrt{Liouville}'s Journ., 2. Ser., Bd. 4, S. 105 (1859).}

If $p$ is a prime, then all the binomial coefficients for the $p$th power

\[
B_h^{(p)} = \dfrac{p(p-1)\ldots(p-h+1)}{1\,.\,2\ldots h},
\]

are, with the exception of $B_0^{(p)}$ and $B_p^{(p)}$, integers divisible by $p$,
%CONT>>
